% ============================================================
% CHAPTER 12 — Project 9: Flutter and Backend: Authentication and CRUD
% ============================================================
\chapter{Project 9: Flutter and Backend: Authentication and CRUD}\index{Flutter!autenticazione}\index{deep link}

\section*{What You'll Build}
A fully working Flutter app that:
\begin{itemize}
  \item Performs OAuth~2.0 login (Google/GitHub) via WebView
  \item Stores tokens securely with \texttt{flutter\_secure\_storage}
  \item Shows the notes list from the API with full CRUD
  \item Handles network errors and expired tokens
  \item Works on Android emulator, iOS, and Chrome
\end{itemize}

\textbf{Estimated time}: 60--90 minutes\\
\textbf{Prerequisite}: Flutter app with static UI (Ch.~11) + Working backend

% ---------------------------------------------------------
\section{OAuth Authentication on Mobile}

On the web, OAuth works with browser redirects. On mobile, the flow is different:

\begin{figure}[H]
\centering
\includegraphics[width=0.8\textwidth]{08-oauth-mobile}
\caption{OAuth flow on mobile}
\end{figure}

\begin{praticabox}[title={Hands-on --- Implement mobile OAuth}]
\begin{lstlisting}[language={}]
1. Add flutter_web_auth_2 to pubspec.yaml
2. In auth_service.dart: open a WebView with the backend
   OAuth URL, intercept the callback scheme "notesapp://",
   extract the token and save it with flutter_secure_storage
3. In the backend: if query ?platform=mobile, redirect to
   notesapp://auth/callback?token=...&refresh=...
\end{lstlisting}
\end{praticabox}

\begin{warnbox}
Passing tokens as query parameters is acceptable for mobile deep links (the \texttt{notesapp://} scheme cannot be intercepted by browsers). \textsc{Do not} ever use this technique on \texttt{https://} URLs.
\end{warnbox}

% ---------------------------------------------------------
\section{Android and iOS Configuration}

\textbf{Android} --- add to \texttt{AndroidManifest.xml}:
\begin{lstlisting}[language=XML]
<intent-filter>
  <action android:name="android.intent.action.VIEW" />
  <category android:name="android.intent.category.DEFAULT" />
  <category android:name="android.intent.category.BROWSABLE" />
  <data android:scheme="notesapp" android:host="auth" />
</intent-filter>
\end{lstlisting}

\textbf{iOS} --- add to \texttt{Info.plist}:
\begin{lstlisting}[language=XML]
<key>CFBundleURLTypes</key>
<array>
  <dict>
    <key>CFBundleURLSchemes</key>
    <array>
      <string>notesapp</string>
    </array>
  </dict>
</array>
\end{lstlisting}

% ---------------------------------------------------------
\section{Secure Token Storage}

\begin{praticabox}[title={Hands-on --- Token service}]
\begin{lstlisting}[language={}]
Create lib/services/token_service.dart that:
1. Uses flutter_secure_storage to save/read/delete
2. Saves: access_token, refresh_token
3. Exposes: saveTokens(), getAccessToken(),
   getRefreshToken(), clearTokens()
4. Android: EncryptedSharedPreferences
5. iOS: Keychain

Then in api_service.dart (Dio):
- Request interceptor: reads token, adds Authorization
- 401 response interceptor: tries refresh, if it fails logout
\end{lstlisting}
\end{praticabox}

\begin{notebox}
\texttt{flutter\_secure\_storage} uses the iOS Keychain and Android EncryptedSharedPreferences --- the operating system's native encrypted storage, much more secure than SharedPreferences.
\end{notebox}

% ---------------------------------------------------------
\section{Connecting the Screens to the Backend}

\begin{praticabox}[title={Hands-on --- Working login}]
\begin{lstlisting}[language={}]
Update login_screen.dart:
- Google button -> authService.loginWithGoogle()
- GitHub button -> authService.loginWithGithub()
- During login: CircularProgressIndicator
- Success: navigate to /dashboard
- Error: SnackBar with message
\end{lstlisting}
\end{praticabox}

\begin{praticabox}[title={Hands-on --- Dashboard with real data}]
\begin{lstlisting}[language={}]
Update dashboard_screen.dart:
1. On mount, call notesProvider to load notes
2. CircularProgressIndicator during loading
3. EmptyState if there are no notes
4. ErrorBanner if there is a network error
5. NoteCard list with real data
6. Pull-to-refresh to reload
7. FAB to create a new note
\end{lstlisting}
\end{praticabox}

\begin{praticabox}[title={Hands-on --- Full CRUD}]
\begin{lstlisting}[language={}]
1. note_form_screen.dart:
   - Form with title and content
   - Validation: title required, min 3 characters
   - Save button: create or update via API
   - If editing, pre-fill the fields

2. note_detail_screen.dart:
   - Show title, content, dates
   - AppBar: edit and delete actions
   - Delete: confirmation dialog

3. dashboard_screen.dart:
   - Tap on NoteCard -> note_detail_screen
   - Swipe-to-delete with confirmation
\end{lstlisting}
\end{praticabox}

% ---------------------------------------------------------
\section{State Management with Riverpod}

The AI generated the providers. Verify this structure:

\begin{lstlisting}[language=Java,title={Riverpod AuthNotifier pattern}]
class AuthNotifier extends StateNotifier<AuthState> {
  final AuthService _authService;
  final TokenService _tokenService;

  AuthNotifier(this._authService, this._tokenService)
      : super(const AuthState()) {
    _checkAuth(); // Auto-login on restart
  }

  Future<void> _checkAuth() async {
    final token = await _tokenService.getAccessToken();
    if (token != null) {
      try {
        final user = await _authService.getCurrentUser();
        state = AuthState(
          status: AuthStatus.authenticated, user: user);
      } catch (_) {
        await _tokenService.clearTokens();
        state = const AuthState(
          status: AuthStatus.unauthenticated);
      }
    } else {
      state = const AuthState(
        status: AuthStatus.unauthenticated);
    }
  }
}
\end{lstlisting}

\begin{tipbox}
Verify that \texttt{\_checkAuth()} is called in the constructor. It is the ``auto-login'' pattern: if there is a saved token, the app logs in automatically without repeating OAuth.
\end{tipbox}

% ---------------------------------------------------------
\section{Testing the App}

\begin{center}
\begin{tabularx}{\textwidth}{l X}
\toprule
\textbf{Feature} & \textbf{Test} \\
\midrule
Google login & Tap $\rightarrow$ WebView $\rightarrow$ Consent $\rightarrow$ Dashboard \\
GitHub login & Tap $\rightarrow$ WebView $\rightarrow$ Consent $\rightarrow$ Dashboard \\
Auto-login & Close and reopen app $\rightarrow$ still authenticated \\
Notes list & Notes from the backend appear \\
Create note & FAB $\rightarrow$ form $\rightarrow$ save $\rightarrow$ note appears \\
Edit note & Tap $\rightarrow$ edit $\rightarrow$ modify $\rightarrow$ save \\
Delete note & Swipe or icon $\rightarrow$ confirm $\rightarrow$ disappears \\
Pull refresh & Pull down $\rightarrow$ list refreshed \\
Network error & Backend off $\rightarrow$ visible error \\
Expired token & Automatic renewal \\
Logout & Tap $\rightarrow$ login $\rightarrow$ tokens cleared \\
\bottomrule
\end{tabularx}
\end{center}

\begin{checkpointbox}
If you can log in, create/edit/delete a note from the mobile app, and the data also appears in the web frontend, the integration is complete.
\end{checkpointbox}

% ---------------------------------------------------------
\section*{Summary}

\begin{center}
\begin{tabularx}{\textwidth}{l X}
\toprule
\textbf{Aspect} & \textbf{Detail} \\
\midrule
Mobile auth & OAuth via WebView + deep link callback \\
Token storage & flutter\_secure\_storage (Keychain/EncryptedSharedPrefs) \\
HTTP & Dio with interceptor for token and refresh \\
State & Riverpod with StateNotifier \\
CRUD & Complete: list, create, edit, delete \\
Auto-login & Saved token $\rightarrow$ automatic login on restart \\
\bottomrule
\end{tabularx}
\end{center}
